% =============================================================================
% Chapter: OLAP & Data Warehousing
% =============================================================================

\chapter{OLAP \& Data Warehousing}

\begin{formulabox}[Key Operations - MEMORIZE!]
\textbf{OLAP Operations:}
\begin{center}
\begin{tabular}{lp{8cm}}
\toprule
\textbf{Operation} & \textbf{Description} \\
\midrule
\textbf{Roll-up} & Aggregation: λεπτομέρεια → σύνοψη (π.χ. day → month) \\
\textbf{Drill-down} & Disaggregation: σύνοψη → λεπτομέρεια (π.χ. year → quarter) \\
\textbf{Slice} & Selection on ONE dimension (π.χ. time = 'Q1') \\
\textbf{Dice} & Selection on MULTIPLE dimensions \\
\textbf{Pivot} & Rotate cube (swap dimensions) \\
\bottomrule
\end{tabular}
\end{center}

\textbf{Schema Types:}
\begin{itemize}
    \item \textbf{Star Schema:} Central fact table + dimension tables (denormalized)
    \item \textbf{Snowflake Schema:} Dimension tables are normalized
\end{itemize}
\end{formulabox}

% =============================================================================
\section{Data Warehouse Concepts}

\subsection{OLTP vs OLAP}

\begin{center}
\begin{tikzpicture}
    \matrix (m) [
        matrix of nodes,
        nodes={
            draw=primary,
            minimum width=4cm,
            minimum height=1.2cm,
            anchor=center,
            font=\small
        },
        column 1/.style={nodes={fill=ntua-blue!20, font=\bfseries}},
        row 1/.style={nodes={fill=secondary!30, font=\bfseries}},
        column sep=-\pgflinewidth,
        row sep=-\pgflinewidth
    ] {
        Aspect & OLTP & OLAP \\
        Purpose & Daily operations & Analysis/Reporting \\
        Data & Current, detailed & Historical, summarized \\
        Queries & Simple, frequent & Complex, ad-hoc \\
        Updates & Many, small & Periodic batch loads \\
        Users & Clerks, customers & Analysts, managers \\
    };
\end{tikzpicture}
\end{center}

\subsection{Data Cube Structure}

\begin{center}
\begin{tikzpicture}[
    cube/.style={draw=primary, thick, fill=secondary!20},
    axis/.style={->, thick}
]
    % Draw cube
    \draw[cube] (0,0,0) -- (3,0,0) -- (3,3,0) -- (0,3,0) -- cycle;
    \draw[cube] (0,0,0) -- (0,0,3) -- (0,3,3) -- (0,3,0);
    \draw[cube] (0,0,0) -- (3,0,0) -- (3,0,3) -- (0,0,3);
    \draw[cube] (3,0,0) -- (3,0,3) -- (3,3,3) -- (3,3,0);
    \draw[cube] (0,3,0) -- (3,3,0) -- (3,3,3) -- (0,3,3);
    \draw[cube] (0,0,3) -- (3,0,3) -- (3,3,3) -- (0,3,3);

    % Axes labels
    \node[anchor=north] at (1.5,-0.3,0) {\textbf{Product}};
    \node[anchor=east] at (-0.3,1.5,0) {\textbf{Time}};
    \node[anchor=south east] at (0,0,4) {\textbf{Location}};

    % Measure
    \node[anchor=west, font=\small] at (4,1.5,0) {Measure: Sales};
\end{tikzpicture}
\end{center}

% =============================================================================
\section{Star vs Snowflake Schema}

\subsection{Star Schema}

\begin{center}
\begin{tikzpicture}[
    fact/.style={draw=accent, fill=accent!20, rounded corners=5pt, minimum width=2.5cm, minimum height=1.5cm},
    dim/.style={draw=secondary, fill=box-blue, rounded corners=5pt, minimum width=2cm, minimum height=1cm}
]
    % Fact table
    \node[fact] (fact) at (0,0) {\textbf{Sales Fact}};

    % Dimension tables
    \node[dim] (time) at (-3,2) {Time};
    \node[dim] (product) at (3,2) {Product};
    \node[dim] (store) at (-3,-2) {Store};
    \node[dim] (customer) at (3,-2) {Customer};

    % Connections
    \draw[thick] (fact) -- (time);
    \draw[thick] (fact) -- (product);
    \draw[thick] (fact) -- (store);
    \draw[thick] (fact) -- (customer);
\end{tikzpicture}
\end{center}

\begin{itemize}
    \item \textbf{Advantages:} Simple queries, fewer joins
    \item \textbf{Disadvantage:} Data redundancy in dimensions
\end{itemize}

\subsection{Snowflake Schema}

\begin{center}
\begin{tikzpicture}[
    fact/.style={draw=accent, fill=accent!20, rounded corners=5pt, minimum width=2cm, minimum height=1cm},
    dim/.style={draw=secondary, fill=box-blue, rounded corners=5pt, minimum width=1.5cm, minimum height=0.8cm, font=\small},
    subdim/.style={draw=success, fill=box-green, rounded corners=3pt, minimum width=1.2cm, minimum height=0.6cm, font=\footnotesize}
]
    % Fact table
    \node[fact] (fact) at (0,0) {\textbf{Sales}};

    % Dimension tables
    \node[dim] (time) at (-2.5,1.5) {Time};
    \node[dim] (product) at (2.5,1.5) {Product};
    \node[dim] (store) at (-2.5,-1.5) {Store};

    % Sub-dimensions (normalized)
    \node[subdim] (month) at (-4.5,2.5) {Month};
    \node[subdim] (year) at (-4.5,1.5) {Year};
    \node[subdim] (category) at (4.5,2.5) {Category};
    \node[subdim] (brand) at (4.5,1.5) {Brand};
    \node[subdim] (city) at (-4.5,-1.5) {City};
    \node[subdim] (region) at (-4.5,-2.5) {Region};

    % Connections
    \draw[thick] (fact) -- (time);
    \draw[thick] (fact) -- (product);
    \draw[thick] (fact) -- (store);

    \draw (time) -- (month);
    \draw (time) -- (year);
    \draw (product) -- (category);
    \draw (product) -- (brand);
    \draw (store) -- (city);
    \draw (city) -- (region);
\end{tikzpicture}
\end{center}

\begin{itemize}
    \item \textbf{Advantage:} Less redundancy, normalized
    \item \textbf{Disadvantage:} More complex queries, more joins
\end{itemize}

% =============================================================================
\section{OLAP Operations (EXAM QUESTIONS!)}

\subsection{Roll-up \& Drill-down}

\begin{center}
\begin{tikzpicture}
    % Roll-up
    \node[draw=secondary, fill=box-blue, rounded corners=8pt, minimum width=5cm, minimum height=3cm] (rollup) at (0,0) {};
    \node[anchor=north, font=\bfseries] at (rollup.north) {Roll-up (Aggregation)};
    \node[anchor=center, align=center, font=\small] at (rollup.center) {
        Day → Month → Quarter → Year\\[0.5em]
        City → State → Country\\[0.5em]
        \textcolor{success}{SUM, COUNT, AVG}
    };

    % Arrow
    \draw[<->, thick, primary] (3,0) -- (5,0);

    % Drill-down
    \node[draw=accent, fill=accent!20, rounded corners=8pt, minimum width=5cm, minimum height=3cm] (drilldown) at (8,0) {};
    \node[anchor=north, font=\bfseries] at (drilldown.north) {Drill-down (Detail)};
    \node[anchor=center, align=center, font=\small] at (drilldown.center) {
        Year → Quarter → Month → Day\\[0.5em]
        Country → State → City\\[0.5em]
        \textcolor{accent}{More granular data}
    };
\end{tikzpicture}
\end{center}

\subsection{Slice \& Dice}

\begin{center}
\begin{tikzpicture}
    % Slice
    \node[draw=success, fill=box-green, rounded corners=8pt, minimum width=5cm, minimum height=2.5cm] (slice) at (0,0) {};
    \node[anchor=north, font=\bfseries] at (slice.north) {Slice};
    \node[anchor=center, align=center, font=\small] at (slice.center) {
        Selection on \textbf{ONE} dimension\\
        e.g., Time = 'Q1 2024'
    };

    % Dice
    \node[draw=warning, fill=warning!20, rounded corners=8pt, minimum width=5cm, minimum height=2.5cm] (dice) at (7,0) {};
    \node[anchor=north, font=\bfseries] at (dice.north) {Dice};
    \node[anchor=center, align=center, font=\small] at (dice.center) {
        Selection on \textbf{MULTIPLE} dimensions\\
        e.g., Time='Q1' AND Product='TV'
    };
\end{tikzpicture}
\end{center}

\subsection{Pivot (Rotate)}

\begin{definitionbox}[Pivot]
Αλλαγή της οπτικής γωνίας του cube - εναλλαγή dimensions στους άξονες.

\textbf{Example:}
\begin{center}
\begin{tabular}{l|cc}
& Q1 & Q2 \\
\hline
TV & 100 & 150 \\
Radio & 50 & 75 \\
\end{tabular}
$\xrightarrow{\text{Pivot}}$
\begin{tabular}{l|cc}
& TV & Radio \\
\hline
Q1 & 100 & 50 \\
Q2 & 150 & 75 \\
\end{tabular}
\end{center}
\end{definitionbox}

% =============================================================================
\section{Solved Exam Problems}

\subsection{2024-25 Exam Questions 1-3 Style}

\begin{examplebox}[OLAP Operations Identification]
\textbf{Sales Data Cube:} Dimensions = (Time, Product, Location)

\textbf{Question 1:} View sales by quarter instead of month.

\textbf{Answer:} \textbf{Roll-up} on Time dimension (month → quarter)

\textbf{Question 2:} View sales for Q1 2024 only.

\textbf{Answer:} \textbf{Slice} (selection on Time = 'Q1 2024')

\textbf{Question 3:} View sales for Q1 2024, in Athens, for Electronics.

\textbf{Answer:} \textbf{Dice} (selection on multiple dimensions)
\end{examplebox}

\subsection{2024-25 Exam Question 15 Style}

\begin{examplebox}[Data Cube Construction]
\textbf{Given:} Fact table with (Product, Store, Time, Sales)

\textbf{Question:} How many cells in a data cube with:
\begin{itemize}
    \item 10 products
    \item 5 stores
    \item 4 quarters
\end{itemize}

\textbf{Answer:}
\begin{itemize}
    \item Base cuboid: $10 \times 5 \times 4 = 200$ cells
    \item With aggregates (ALL):
    \begin{itemize}
        \item Product × Store × Time: 200
        \item Product × Store × ALL: $10 \times 5 = 50$
        \item Product × ALL × Time: $10 \times 4 = 40$
        \item ALL × Store × Time: $5 \times 4 = 20$
        \item Product × ALL × ALL: 10
        \item ALL × Store × ALL: 5
        \item ALL × ALL × Time: 4
        \item ALL × ALL × ALL: 1
    \end{itemize}
    \item Total: $200 + 50 + 40 + 20 + 10 + 5 + 4 + 1 = 330$ cells
\end{itemize}

\textbf{General Formula:} For dimensions with $d_1, d_2, ..., d_n$ values:
\[
\text{Total cells} = (d_1 + 1)(d_2 + 1)...(d_n + 1)
\]
Here: $(10+1)(5+1)(4+1) = 11 \times 6 \times 5 = 330$
\end{examplebox}

% =============================================================================
\section{Common Mistakes}

\begin{warningbox}[Συχνά Λάθη στις Εξετάσεις]
\begin{enumerate}
    \item \textbf{Roll-up vs Drill-down confusion}:
    \begin{itemize}
        \item Roll-up: Less detail, MORE aggregation (day→month)
        \item Drill-down: More detail, LESS aggregation (month→day)
    \end{itemize}

    \item \textbf{Slice vs Dice}:
    \begin{itemize}
        \item Slice: ONE dimension selection
        \item Dice: MULTIPLE dimension selection
    \end{itemize}

    \item \textbf{Star vs Snowflake}:
    \begin{itemize}
        \item Star: Denormalized dimensions (faster queries)
        \item Snowflake: Normalized dimensions (less redundancy)
    \end{itemize}

    \item \textbf{Cube cell count}:
    \begin{itemize}
        \item Don't forget the ALL aggregations (+1 for each dimension)
    \end{itemize}
\end{enumerate}
\end{warningbox}

% =============================================================================
\section{Quick Reference}

\begin{center}
\begin{tikzpicture}
    \node[draw=ntua-blue, fill=ntua-blue!10, rounded corners=5pt, minimum width=14cm, minimum height=7cm] (main) {};

    \node[anchor=north west, font=\bfseries\Large\color{ntua-blue}] at ([xshift=5pt, yshift=-5pt]main.north west) {OLAP Operations Summary};

    \node[anchor=north west, align=left, font=\small] at ([xshift=10pt, yshift=-35pt]main.north west) {
        \textbf{OLAP Operations:}\\
        • \textbf{Roll-up}: Aggregate (less detail): day→month→year\\
        • \textbf{Drill-down}: Disaggregate (more detail): year→month→day\\
        • \textbf{Slice}: Filter ONE dimension\\
        • \textbf{Dice}: Filter MULTIPLE dimensions\\
        • \textbf{Pivot}: Rotate axes\\[5pt]
        \textbf{Schemas:}\\
        • Star: Fact table + denormalized dimensions (simple, fast)\\
        • Snowflake: Normalized dimensions (complex, less redundant)
    };

    \node[draw=secondary, fill=box-blue, rounded corners=5pt, minimum width=5cm, minimum height=1.8cm]
        at ([xshift=-4cm, yshift=1.5cm]main.south) {};
    \node[align=center, font=\footnotesize] at ([xshift=-4cm, yshift=1.5cm]main.south) {
        \textbf{Cube Cells:}\\
        $(d_1+1)(d_2+1)...(d_n+1)$\\
        (+1 for ALL aggregation)
    };

    \node[draw=accent, fill=box-red, rounded corners=5pt, minimum width=5cm, minimum height=1.8cm]
        at ([xshift=3cm, yshift=1.5cm]main.south) {};
    \node[align=center, font=\footnotesize] at ([xshift=3cm, yshift=1.5cm]main.south) {
        \textbf{Remember:}\\
        Roll-up = UP the hierarchy\\
        Drill-down = DOWN to detail
    };
\end{tikzpicture}
\end{center}
